\section{Modern C++ Features}

C++ has evolved significantly since C++11, with each standard bringing features that make the language safer, more expressive, and easier to use. This section covers key features from C++11 through C++20 that every modern C++ developer should know.

\subsection{Type Inference with \texttt{auto}}

The \texttt{auto} keyword lets the compiler deduce the type of a variable from its initializer:

\begin{lstlisting}
	// Without auto
	std::vector<int>::iterator it = vec.begin();
	std::map<std::string, std::vector<int>>::const_iterator mit = myMap.find("key");
	
	// With auto - much cleaner!
	auto it = vec.begin();
	auto mit = myMap.find("key");
	
	// Works with any initializer
	auto x = 42;           // int
	auto y = 3.14;         // double
	auto s = "hello"s;     // std::string (with string literal suffix)
	auto ptr = std::make_unique<Widget>();  // std::unique_ptr<Widget>
\end{lstlisting}

\begin{remark}
	Use \texttt{auto} when the type is obvious from context or overly verbose. Avoid it when the type provides important documentation or when it might be ambiguous.
\end{remark}

\subsection{Range-Based For Loops}

Iterate over containers without dealing with iterators:

\begin{lstlisting}
	std::vector<int> numbers = {1, 2, 3, 4, 5};
	
	// Old style
	for (std::vector<int>::iterator it = numbers.begin(); it != numbers.end(); ++it) {
		std::cout << *it << " ";
	}
	
	// Range-based for loop
	for (int n : numbers) {
		std::cout << n << " ";
	}
	
	// With reference to modify elements
	for (int& n : numbers) {
		n *= 2;
	}
	
	// With auto
	for (auto& n : numbers) {
		n *= 2;
	}
	
	// With const reference (read-only, no copies)
	for (const auto& n : numbers) {
		std::cout << n << " ";
	}
\end{lstlisting}

\subsection{Uniform Initialization}

C++11 introduced brace initialization, which works consistently across all types:

\begin{lstlisting}
	// All of these use brace initialization
	int x{42};
	double d{3.14};
	std::string s{"hello"};
	std::vector<int> v{1, 2, 3, 4, 5};
	std::map<std::string, int> m{{"alice", 25}, {"bob", 30}};
	
	// Works with custom types too
	struct Point { int x, y; };
	Point p{10, 20};
	
	// Prevents narrowing conversions
	int a{3.14};    // ERROR: narrowing from double to int
	int b = 3.14;   // OK but loses precision (b = 3)
\end{lstlisting}

\subsection{nullptr}

\texttt{nullptr} is a type-safe null pointer constant that replaces \texttt{NULL}:

\begin{lstlisting}
	// Old style (problematic)
	void foo(int);
	void foo(int*);
	foo(NULL);      // Ambiguous! NULL is typically defined as 0
	
	// Modern style
	foo(nullptr);   // Unambiguous - calls foo(int*)
	
	int* p = nullptr;
	if (p == nullptr) {
		// ...
	}
\end{lstlisting}

\subsection{constexpr}

\texttt{constexpr} indicates that a value or function can be evaluated at compile time:

\begin{lstlisting}
	// Compile-time constants
	constexpr int MAX_SIZE = 100;
	constexpr double PI = 3.14159265359;
	
	// Compile-time functions
	constexpr int factorial(int n) {
		return (n <= 1) ? 1 : n * factorial(n - 1);
	}
	
	constexpr int fact5 = factorial(5);  // Computed at compile time!
	
	// Can be used where compile-time constants are required
	int arr[factorial(5)];  // Array of 120 elements
\end{lstlisting}

\begin{remark}
	C++14 relaxed \texttt{constexpr} function restrictions (allowing loops, local variables, etc.), and C++20 expanded it further with \texttt{consteval} for functions that \textit{must} be evaluated at compile time.
\end{remark}

\subsection{Structured Bindings (C++17)}

Structured bindings allow you to decompose objects into their components:

\begin{lstlisting}
	// With pairs
	std::pair<int, std::string> getPerson() {
		return {25, "Alice"};
	}
	
	auto [age, name] = getPerson();
	std::cout << name << " is " << age << " years old\n";
	
	// With maps
	std::map<std::string, int> scores = {{"Alice", 95}, {"Bob", 87}};
	
	for (const auto& [name, score] : scores) {
		std::cout << name << ": " << score << "\n";
	}
	
	// With structs
	struct Point { int x, y, z; };
	Point p{1, 2, 3};
	auto [x, y, z] = p;
	
	// With arrays
	int arr[] = {1, 2, 3};
	auto [a, b, c] = arr;
\end{lstlisting}

\subsection{std::optional (C++17)}

\texttt{std::optional} represents a value that may or may not be present:

\begin{lstlisting}
	#include <optional>
	
	std::optional<int> findIndex(const std::vector<int>& vec, int value) {
		for (size_t i = 0; i < vec.size(); ++i) {
			if (vec[i] == value) {
				return i;  // Return the index
			}
		}
		return std::nullopt;  // No value found
	}
	
	auto result = findIndex(myVec, 42);
	
	if (result) {  // or result.has_value()
		std::cout << "Found at index: " << *result << "\n";  // or result.value()
	} else {
		std::cout << "Not found\n";
	}
	
	// With value_or for defaults
	int index = findIndex(myVec, 42).value_or(-1);
\end{lstlisting}

\subsection{std::variant (C++17)}

\texttt{std::variant} is a type-safe union that can hold one of several specified types:

\begin{lstlisting}
	#include <variant>
	
	std::variant<int, double, std::string> value;
	
	value = 42;                    // Holds int
	value = 3.14;                  // Now holds double
	value = "hello";               // Now holds string
	
	// Check what it holds
	if (std::holds_alternative<int>(value)) {
		std::cout << "It's an int: " << std::get<int>(value) << "\n";
	}
	
	// Visit pattern - handle all alternatives
	std::visit([](auto&& arg) {
		std::cout << arg << "\n";
	}, value);
	
	// Visitor with overloaded lambdas
	std::visit(overloaded{
		[](int i) { std::cout << "int: " << i << "\n"; },
		[](double d) { std::cout << "double: " << d << "\n"; },
		[](const std::string& s) { std::cout << "string: " << s << "\n"; }
	}, value);
\end{lstlisting}

\subsection{if/switch with Initializers (C++17)}

You can now include an initialization statement in \texttt{if} and \texttt{switch}:

\begin{lstlisting}
	// Old style
	auto it = myMap.find(key);
	if (it != myMap.end()) {
		// use it->second
	}
	// it is still in scope here
	
	// C++17 style
	if (auto it = myMap.find(key); it != myMap.end()) {
		// use it->second
	}
	// it is out of scope here - cleaner!
	
	// Works with switch too
	switch (auto ch = getChar(); ch) {
		case 'a': /* ... */ break;
		case 'b': /* ... */ break;
		default: /* ... */ break;
	}
\end{lstlisting}

\subsection{String View (C++17)}

\texttt{std::string\_view} is a non-owning view into a string, avoiding copies:

\begin{lstlisting}
	#include <string_view>
	
	// Old style - may copy the string
	void processString(const std::string& s);
	
	// Better - no copies, works with string literals and std::string
	void processString(std::string_view s) {
		std::cout << s.substr(0, 5) << "\n";  // Also returns string_view
	}
	
	processString("hello world");           // No allocation
	processString(std::string("hello"));    // No copy of the string data
	
	// Creating string_view
	std::string_view sv1 = "hello";
	std::string str = "world";
	std::string_view sv2 = str;
\end{lstlisting}

\begin{remark}
	Be careful with \texttt{string\_view} lifetime—it doesn't own the data, so make sure the underlying string outlives the view.
\end{remark}

\subsection{Concepts (C++20)}

Concepts constrain template parameters, providing clear error messages and self-documenting code:

\begin{lstlisting}
	#include <concepts>
	
	// Define a concept
	template <typename T>
	concept Numeric = std::integral<T> || std::floating_point<T>;
	
	// Use the concept
	template <Numeric T>
	T add(T a, T b) {
		return a + b;
	}
	
	// Alternative syntax
	template <typename T>
	requires Numeric<T>
	T multiply(T a, T b) {
		return a * b;
	}
	
	// Abbreviated syntax
	auto divide(Numeric auto a, Numeric auto b) {
		return a / b;
	}
	
	add(5, 3);          // OK
	add(3.14, 2.71);    // OK
	add("hello", "world");  // Error: strings don't satisfy Numeric
\end{lstlisting}

\subsection{Ranges (C++20)}

Ranges provide a composable way to work with sequences:

\begin{lstlisting}
	#include <ranges>
	#include <algorithm>
	
	std::vector<int> numbers = {1, 2, 3, 4, 5, 6, 7, 8, 9, 10};
	
	// Filter and transform with views (lazy evaluation)
	auto result = numbers 
	| std::views::filter([](int n) { return n % 2 == 0; })  // Keep evens
	| std::views::transform([](int n) { return n * n; });   // Square them
	
	for (int n : result) {
		std::cout << n << " ";  // 4 16 36 64 100
	}
	
	// Take first N elements
	for (int n : numbers | std::views::take(3)) {
		std::cout << n << " ";  // 1 2 3
	}
	
	// Reverse
	for (int n : numbers | std::views::reverse) {
		std::cout << n << " ";  // 10 9 8 7 6 5 4 3 2 1
	}
\end{lstlisting}

\subsection{Three-Way Comparison (C++20)}

The spaceship operator (\texttt{<=>}) simplifies comparison operations:

\begin{lstlisting}
	#include <compare>
	
	struct Point {
		int x, y;
		
		// One operator generates all six comparison operators!
		auto operator<=>(const Point&) const = default;
	};
	
	Point p1{1, 2}, p2{1, 3};
	
	if (p1 < p2) { /* ... */ }   // Works!
	if (p1 == p2) { /* ... */ }  // Works!
	if (p1 >= p2) { /* ... */ }  // Works!
	
	// Custom implementation
	struct Version {
		int major, minor, patch;
		
		auto operator<=>(const Version& other) const {
			if (auto cmp = major <=> other.major; cmp != 0) return cmp;
			if (auto cmp = minor <=> other.minor; cmp != 0) return cmp;
			return patch <=> other.patch;
		}
		
		bool operator==(const Version&) const = default;
	};
\end{lstlisting}

\subsection{Summary: Which Standard to Use}

\begin{itemize}
	\item \textbf{C++11}: The minimum for modern C++. Includes move semantics, \texttt{auto}, range-for, lambdas, smart pointers.
	\item \textbf{C++14}: Minor improvements to C++11. Better lambdas, relaxed \texttt{constexpr}.
	\item \textbf{C++17}: Major quality-of-life improvements. Structured bindings,
		\texttt{optional}, \texttt{variant}, \texttt{string\_view}, if-initializers.
	\item \textbf{C++20}: Significant new features. Concepts, ranges, coroutines, modules, three-way comparison.
	\item \textbf{C++23}: Latest standard with further refinements.
\end{itemize}

For new projects, target at least C++17. Use C++20 if your compiler and team support it.